\chapter{Conclusion}
\label{Chapter5}

A general nonlinear framework for the solution of rotation-free solid membranes is presented. The theory in Section \ref{Chapter2} details
 the solution of general elliptic partial differential equations using the Finite Element Method. The function spaces of the test
 and trial functions are introduced for the well-posedness of the weak formulation necessary for a meaningful solution of the
 boundary value problem. Dirichlet, Neumann, and mixed boundary conditions are considered and the isoparametric parametric concept
 along with numerical quadrature is used to simplify the implementation of the solution on MATLAB. Section \ref{Chapter3} develops a
 nonlinear framework for the solution of problems involving large deformations of beams in a Euclidean basis. Discretisation of
 general kinematical quantities, balance laws, and constitution is introduced. The resulting nonlinear system of equations is solved using the Newton-Raphson method. 
 This framework is validated using simple numerical examples while comprehensively covering the load and displacement driven boundary conditions. 
 After a thorough development of the foundation, Section \ref{Chapter4} extends the previous chapters to develop the framework for curved structures.
 Based on the differential geometry of curved surfaces, the resulting framework is general, accounts for large deformations, and general material laws.
 The use of curvilinear coordinates further eases computation as the resulting formulation only uses three degrees of freedom, avoids the use
 of a local cartesian coordinate system, the transformation of derivatives, and allows the splitting of the strong and weak forms
 into the in-plane and out-of-plane parts allowing an easy extension of the framework to liquid membranes. The solution of the final nonlinear system
 is conducted using the theory developed in the previous chapters with the volume constraints being included in the system using
 the Lagrange multiplier method. This framework is validated through the consideration of the inflation of a square sheet. It is expected
 that the MATLAB implementation of the formuation will be able to reproduce the results of \cite{sauer2014} in the future. 
